% !TEX root = ../../main_docker.tex
% Schlanke Images

\section{Mehrstufige Builds und schlanke Images}
\label{sec:multistage}

\subsection{Dangling Images}
Images ohne Namen (engl.\ „baumelnd“, wie der dangling pointer in C).
Zwei Entstehungswege:
\begin{enumerate}
\item \R{docker build .} ohne \R{-t} \ra das Image bekommt nie einen Namen
\item ein neueres Image gleichen Namens \T{nimmt dem älteren den Namen weg}
      -- beide bleiben erhalten, denn es könnte noch einen Container der
      alten Version geben
\end{enumerate}
\begin{ztab}{@{}lX@{}}
\R{docker images -a} & zeigt sie als \R{<untagged>} \\
\R{docker images --no-trunc} & zeigt \R{<none>:<none>} \\
\R{docker images --filter dangling=true} & nur die namenlosen \\
\R{docker image prune} & alle dangling Images löschen \\
\R{docker image prune -a} & \emph{alle} Images ohne zugeordneten Container \\
\R{docker tag <alt> <neu>} & nachträglich (um)benennen \\
\end{ztab}
\merkzeile Ein versehentlich überschriebener Tag ist also kein
Datenverlust: Image-ID aus \R{docker images -a} holen und mit
\R{docker tag} wieder benennen -- solange man nicht vorher \R{prune} laufen lässt.

\subsection{Multi-stage Builds}
\label{rez:multistage}
\begin{formel}[Die Grundidee]
Das übersetzte Programm läuft auch in einem Image \emph{ohne} SDK. Also:
in Stufe 1 mit der Entwicklungsumgebung bauen, in Stufe 2 nur das
Ergebnis in ein schlankes Parent-Image kopieren. \T{Mehrere \R{FROM} im
selben Dockerfile $=$ mehrere Stufen; nur die letzte wird zum Resultat.}
\end{formel}
\begin{lstlisting}
FROM gcc as builder                  # Stufe benennen
WORKDIR /myapp
COPY hello-world.c /myapp
RUN gcc -o hello-world hello-world.c

FROM debian                          # neue Stufe = neues Parent-Image
WORKDIR /myapp
COPY --from=builder /myapp/hello-world /myapp
CMD ["/myapp/hello-world"]
\end{lstlisting}
\begin{ztab}{@{}lX@{}}
\R{FROM <image> as <name>} & benennt die Stufe \\
\R{COPY --from=<name> <src> <dest>} & kopiert zwischen Stufen \\
\R{docker build . --target <name>} & nur bis zu dieser Stufe bauen \ra das
  Debug-Werkzeug, wenn ein mehrstufiger Build abbricht und gar kein Image
  entsteht \\
\end{ztab}

\begin{bsp}[Was die Wahl des Parent-Images ausmacht]
Dasselbe \R{hello-world.c}, vier Varianten:
\begin{tabular}{@{}lr@{}}
\R{my-hello-world} (einstufig) & 1.18 GB \\
\R{multistage\_debian} & 118 MB \\
\R{multistage\_alpine} & 8.18 MB \\
\R{multistage\_scratch} & 721 kB \\
\end{tabular}\\[2pt]
Faktor \textbf{$\sim$160} zwischen scratch und debian, Faktor
$\sim$1600 zum einstufigen Build.
\end{bsp}

\subsection{Parent-Image wählen: Funktionalität gegen Grösse}
\begin{ztab}{@{}lrX@{}}
\R{scratch} & 0 B & leeres Image; das Programm muss \emph{alle}
  Abhängigkeiten mitbringen, keine Linux-Kommandos \\
\R{alpine} & $\sim$14 MB & busybox, Shell \R{sh}, \T{musl libc},
  Paketmanager \R{apk}; wichtige Befehle sind Links auf \R{/bin/busybox};
  kein \R{bash}-Komfort (nachrüstbar mit \R{apk add bash}) \\
\R{debian} & $\sim$200 MB & vollständiges Linux: GNU-Werkzeuge,
  \R{bash}, \T{GNU libc}, Paketmanager \R{apt} \\
\end{ztab}
\SW{Was die Anwendung braucht, bestimmt das Parent-Image}
Shell-Skripte \ra Linux mit \R{sh}/\R{bash} $\cdot$ Nachinstallieren
\ra Paketverwaltung $\cdot$ Laufzeitumgebung \ra offizielles Image
(Runtime, nicht SDK) $\cdot$ Drittanbieter-Bibliotheken \ra kompatibel zur
Runtime $\cdot$ im Betrieb debugbar \ra volles Linux mit \R{nano}, \R{ps}.
\merkzeile Der \emph{letzte} Build-Schritt bestimmt Funktionalität und
Grösse des Ergebnisses.

\subsection{Statisch gegen dynamisch binden, und die libc}
\label{sec:linken}
\SW{Compiler und Linker} Der Compiler macht aus jeder \R{.c} eine
\R{.o}-Datei mit den Einträgen „ich brauche X“ und „ich liefere Y“. Der
Linker verbindet jedes „ich brauche X“ mit genau einem „ich liefere X“ --
aus \R{.o}-Dateien und Bibliotheken (\R{.a}, \R{.so}).

\begin{ztab}{@{}lX@{}}
\T{statisch} & \R{gcc -static -o out code.c} -- die benötigten
  Bibliotheksfunktionen werden \emph{ins Binary kopiert}. Grösser, aber in
  sich vollständig. \emph{Die Schnecke kann nicht aus ihrem Haus.} \\
\T{dynamisch} & \R{gcc -o out code.c} -- im Binary steht nur eine
  Referenz auf die shared library; das OS bindet sie beim Programmstart.
  Kleiner, aber abhängig von den installierten Bibliotheken \emph{und
  deren Version}. \emph{Der Einsiedlerkrebs kann das Haus wechseln.} \\
\end{ztab}

\begin{bsp}[Dasselbe Programm, beide Varianten (arm64)]
\begin{tabularx}{\linewidth}{@{}Xr@{}}
dynamisch gebunden, Binary & 9.0 kB \\
\quad dazu \R{libc.so.6} $\to$ \R{libc-2.31.so} im Image & 1.4 MB \\
statisch gebunden, Binary & 706 kB \\
\end{tabularx}\\[2pt]
Geprüft mit \R{ldd hello-world}: dynamisch listet es die shared objects,
statisch antwortet es \R{not a dynamic executable}. Statisch werden
\emph{nur die benötigten} Funktionen kopiert -- nicht die ganzen 1.4 MB.
\end{bsp}

\SW{Pro dynamisch} kleinere Programme, eine Bibliothek für alle $\cdot$
Bibliothek aktualisierbar ohne Neuübersetzung und Auslieferung $\cdot$
effizientes Caching.
\SW{Kontra dynamisch} braucht OS-Unterstützung \ra nicht auf
Mikrocontrollern $\cdot$ Probleme bei veralteten Versionen $\cdot$ läuft
gar nicht, wenn die Bibliothek fehlt.

\SW{libc} Die C-Standardbibliothek: häufig benötigte Grundfunktionen
(Mathematik, Zeichenketten) \emph{und} die User-Space-API zum Kernel
(\R{open}, \R{read}, \R{write}, \R{malloc}, \R{printf}). Auch
Nicht-C-Programme benutzen sie indirekt -- \R{ldd /usr/bin/python} zeigt
\R{libc.so.6}.
\begin{ztab}{@{}lX@{}}
\T{GNU glibc} & in fast allen Distributionen; Ziel: Leistung und
  Portabilität statt geringer Grösse \\
\T{musl libc} & in Alpine Linux; Ziel: Korrektheit und geringe Grösse
  statt Leistung \\
\end{ztab}

\begin{falle}[Warum das dynamisch gebundene Binary auf Alpine nicht lief]
Der \R{gcc} aus dem \R{gcc}-Image bindet gegen \T{glibc}. Alpine bringt
aber \T{musl} mit und kann ein dynamisch gegen glibc gebundenes Binary
nicht starten. Deshalb stehen in \R{multistage\_alpine} und
\R{multistage\_scratch} beide mit \R{gcc -static}.\\
\T{Lösung} einen \R{gcc} verwenden, der gegen musl bindet -- den gibt es
als Alpine-Paket:
\R{RUN apk update \&\& apk add gcc musl-dev}, erste Stufe dann
\R{FROM alpine as builder}.\\
Ergebnis: Binary 70.9 kB, dynamisch gegen
\R{/lib/ld-musl-aarch64.so.1} (642 kB) gebunden. Image 7.53 MB
gegenüber 8.18 MB statisch -- das Alpine-Basis-Image allein sind 7.46 MB.
\end{falle}

\merkzeile Nicht prüfungsrelevant (Folie 362, aus Zeitgründen gestrichen):
GPIO mit Raspi/C\#, Portainer, Docker Networking.

\subsection{Die drei Aufgabenmuster}
\label{rez:muster}
Dieselbe Aufgabe („C\#-Quellen übersetzen und ausführen“) in drei
Varianten -- das ist das Schema, nach dem geprüft wird.
\begin{rezept}[A1: interaktiv, ohne Dockerfile]
Host-Volume einhängen, im Container übersetzen, Ergebnis landet dank
Volume direkt auf dem Host.
\begin{lstlisting}
$ docker run -it --rm -v $(pwd):/work \
    mcr.microsoft.com/dotnet/sdk:7.0 bash
# cd /work
# dotnet publish -c Release -o ./out ./src
# exit
\end{lstlisting}
\merkzeile \R{\$} $=$ Host, \R{\#} $=$ Container. Die Aufgabe verlangt
ausdrücklich \T{Host-Volume statt \R{docker cp}} -- das Volume schreibt
sofort und in beide Richtungen, \R{cp} wäre ein zusätzlicher Schritt
nach dem Lauf.
\end{rezept}
\begin{rezept}[A2: single-stage Dockerfile]
\begin{lstlisting}
FROM mcr.microsoft.com/dotnet/sdk:7.0
WORKDIR /app
COPY . .
RUN dotnet publish -c Release -o out ./src
CMD ["dotnet", "out/Homework.dll"]
\end{lstlisting}
\R{\$ docker build . -t hw} dann \R{\$ docker run --rm hw}.
\T{Nachteil:} das SDK bleibt im Image -- über 1\,GB für ein Programm von
wenigen hundert Kilobyte.
\end{rezept}
\begin{rezept}[A3: multi-stage Dockerfile]
\begin{lstlisting}
FROM mcr.microsoft.com/dotnet/sdk:7.0 as builder
WORKDIR /app
COPY . .
RUN dotnet publish -c Release -o out ./src

FROM mcr.microsoft.com/dotnet/runtime:7.0
WORKDIR /app
COPY --from=builder /app/out /app/out
CMD ["dotnet", "out/Homework.dll"]
\end{lstlisting}
\merkzeile Der Unterschied ist \emph{nur} das Parent-Image der zweiten
Stufe: \T{\R{runtime} statt \R{sdk}}. Zum Ausführen braucht es keinen
Compiler. Bei den offiziellen Images ist das die Standardpaarung --
\R{sdk} zum Bauen, \R{runtime} zum Laufen.
\end{rezept}
\begin{falle}[Diagnose-Variable]
In manchen Vorlagen steht \R{ENV COMPlus\_EnableDiagnostics=0}. Das
schaltet die .NET-Diagnose ab, die sonst im Container über Named Pipes
stolpert. Solche \R{ENV}-Zeilen erzeugen \emph{keinen} Layer -- sie sind
reine Metadaten.
\end{falle}
